%
%
\TODOc{Revise conclusion sections fully - later}
We present a separable volumetric transfer for models with different horizontal meshes and vertical coordinates. The method combines a consistent horizontal remap with a column-wise conservative vertical projection that reproduces \ROMS layer integrals. Because the vertical operator is constrained to preserve these integrals, the composite method is first-order accurate. In the one-dimensional tests, a non-conservative point-value interpolant reproduced smooth profiles with convergence between second and third order, but renders the method non-conservative. 

For practical applications, a \ROMS{}-$\rightarrow$-\MPAS{} resolution ratio of approximately \(\hRatio \approx 0.2\) provides a useful initial target. This value should be tested through a short refinement sequence for each domain rather than adopted as a universal threshold. The case studies also show that remapping alone is insufficient to produce a viable regional configuration. The Bay of Bengal simulation required bathymetric smoothing, while Chesapeake Bay required native estuarine geometry and blended fields in addition to the \MPAS forcing conditions.

Two additional tests would most directly strengthen the verification. First, comparing real fields over their common support on the finest available grid pair would determine whether the apparent error saturation reflects the transfer itself or changing coverage between resolutions. Second, isolated tests of the vertical remapper on synthetic columns would establish its convergence behavior independently of horizontal interpolation and changing wet-domain support. These tests should include cases in which the ROMS water column is contained within, matches, or exceeds the \MPAS depth range. Further evaluation could include forward-transfer tests for physical tracers against surrogate reference solutions, sensitivity experiments using multiple \MPAS meshes, and repeated timing measurements. Together, these tests would provide a stronger basis for evaluating online one-way coupling.

The same framework can support two-way coupling. The roundtrip analysis already requires reverse horizontal and vertical operators, \(W_{R\to M}\) and \(\mathcal{V}_{R\to M}\), in addition to the forward maps. The tensor-product formulation can therefore be applied with ROMS as the source and \MPAS as the target, transferring regional information back to the global mesh.
%What the round trip does not do is validate that reverse leg on its own, because errors incurred on the forward leg can be undone on the return, and the composition is therefore a weaker test than either direction taken separately. 
The numerical considerations differ in the upscaling direction. Fine-to-coarse transfer aggregates many \ROMS cells into a single \MPAS cell, rather than distributing one coarse \MPAS value across many fine \ROMS cells. Consequently, mesh imprinting is not the principal concern, and high-order conservative \(L_2\)-minimization  horizontal transfer can be used without producing the staircase artifacts that motivated use of quasi-interpolatory maps like the normalized bilinear interpolation for downscaling. 

